\documentclass[12pt]{article}
\usepackage{amsmath,amssymb,amsthm}
\usepackage[margin=1in]{geometry}

\newtheorem{theorem}{Theorem}[section]
\newtheorem{lemma}[theorem]{Lemma}
\theoremstyle{definition}
\newtheorem{definition}[theorem]{Definition}
\newtheorem{exercise}[theorem]{Exercise}

\title{Principles of Mathematical Analysis: Chapter 2 Exercises}
\date{}

\begin{document}
\maketitle
\noindent The following solutions emphasize the main argument while retaining the details needed to justify each conclusion.

Prove that the empty set is a subset of every set.

\smallskip
\noindent\textbf{Solution.}\quad

Denote $\varnothing$ by the empty set. Suppose $\varnothing\not\subset E$ for some set $E$, then there exists an element $x\in\varnothing$ such that $x\notin E$. But this violates the definition that $\varnothing$ contains nothing, which is a contradiction.



\bigskip\noindent\hrule\bigskip

A complex number $z$ is said to be \textit{algebraic} if there are integers $a_0,\ldots,a_n$, not all zero, such that
\begin{align*}
a_0z^n+a_1z^{n-1}+\cdots+a_{n-1}z+a_n=0
\end{align*}
Prove that the set of all algebraic numbers is countable. \textit{Hint:} For every positive integer $N$ there are only finitely many equations with
\begin{align*}
n+\left\vert a_0\right\vert+\left\vert a_1\right\vert+\cdots+\left\vert a_n\right\vert=N
\end{align*}

\smallskip
\noindent\textbf{Solution.}\quad

By using the hint, for each $N\in J$, consider the finite set
\begin{align*}
S_N=\left\{(n,a_0,a_1,\ldots,a_n):
\begin{array}{l}
a_0,a_1,\ldots,a_n\mbox{ not all zero, and}\\
n+\left\vert a_0\right\vert+\left\vert a_1\right\vert+\cdots+\left\vert a_n\right\vert=N
\end{array}\right\}
\end{align*}
Moreover, for each $(n,a_0,a_1,\ldots,a_n)\in S_N$, we have the equation
\begin{align*}
a_0z^n+a_1z^{n-1}+\cdots+a_{n-1}z+a_n=0
\end{align*}
Since $a_0,a_1,\ldots,a_n$ not all zero, there are at most $n$ algebraic numbers satisfying the above equation. Let
\begin{align*}
T_N=\bigcup_{(n,a_0,a_1,\ldots,a_n)\in S_N}\{z:a_0z^n+a_1z^{n-1}+\cdots+a_{n-1}z+a_n=0\}
\end{align*}
By the above argument, $T_N$ is a finite set. By the corollary to Theorem 2.12, the set of all algebraic numbers, $\bigcup_{N\in J}T_N$, is at most countable. On the other hand, since every rational number is algebraic, for given $\frac{p}{q}\in Q$, where $p$ and $q$ are integers and $q>0$, we have $q\left(\frac{p}{q}\right)-p=0$. Hence the set of all algebraic numbers is countable.



\bigskip\noindent\hrule\bigskip

Prove that there exist real numbers which are not algebraic.

\smallskip
\noindent\textbf{Solution.}\quad

By the remark of Theorem 2.14, the set of all real numbers is uncountable. So if every real number is algebraic, then by Exercise 2, it is also countable. Hence a contradiction occurs.



\bigskip\noindent\hrule\bigskip

Is the set of all irrational real numbers countable?

\smallskip
\noindent\textbf{Solution.}\quad

No. If the set of all irrational real numbers $R-Q$ is countable, then by Theorem 2.13, the set of all real numbers $R=Q\cup(R-Q)$ is countable. This gives a contradiction.



\bigskip\noindent\hrule\bigskip

Construct a bounded set of real numbers with exactly three limit points.

\smallskip
\noindent\textbf{Solution.}\quad

For example, let $E=\left\{1+\frac{1}{n}:n\in J\right\}\cup\left\{3+\frac{1}{n}:n\in J\right\}\cup\left\{5+\frac{1}{n}:n\in J\right\}$.



\bigskip\noindent\hrule\bigskip

Let $E'$ be the set of all limit points of a set $E$. Prove that $E'$ is closed. Prove that $E$ and $\bar{E}$ have the same limit points. (Recall that $\bar{E}=E\cup E'$.) Do $E$ and $E'$ always have the same limit points?

\smallskip
\noindent\textbf{Solution.}\quad

(i) Let $p$ be a limit point of $E'$. Then given $r>0$, there exists a point $q\ne p$ with $d(p,q)<r$, such that $q\in E'$. Let $d(p,q)=r-h$, and $m=\min\{h,r-h\}>0$. Then there exists a point $s\ne q$ with $d(q,s)<m$, such that $s\in E$. It follows that
\begin{align*}
d(p,s)
\le d(p,q)+d(q,s)
<r-h+m
\le r-h+h
=r
\end{align*}
On the other hand,
\begin{align*}
d(p,s)
\ge d(p,q)-d(q,s)
> r-h-m
\ge 0
\end{align*}
i.e., $s\ne p$. Hence $p\in E'$ and then $E'$ is closed.
(ii) Let $\bar{E}'$ be the set of all limit points of $\bar{E}$. If $p\in E'$, then $p$ is a limit point of $E\subset\bar{E}$, which is clear that $p\in\bar{E}'$. Conversely, if $p\in\bar{E}'$, then for any $r>0$ there exists a point $q\ne p$ with $d(p,q)<r$, such that $q\in\bar{E}=E\cup E'$. If $q\in E$, then the claim follows. If $q\in E'$, then by the previous result, we can find a point $s\ne p$ with $d(p,s)<r$ such that $s\in E$. Thus $p\in E'$.
(iii) The answer is no. For example, as we proposed in Exercise 5, the set
\begin{align*}
E=\left\{1+\frac{1}{n}:n\in J\right\}\cup\left\{3+\frac{1}{n}:n\in J\right\}\cup\left\{5+\frac{1}{n}:n\in J\right\}
\end{align*}
has only three limit points, whereas $E'=\{1,3,5\}$ has none.



\bigskip\noindent\hrule\bigskip

Let $A_1,A_2,A_3,\ldots$ be subsets of a metric space.
(a) If $B_n=\bigcup_{i=1}^nA_i$, prove that $\bar{B}_n=\bigcup_{i=1}^n\bar{A}_i$, for $n=1,2,3,\ldots$.
(b) If $B=\bigcup_{i=1}^\infty A_i$, prove that $\bar{B}\supset\bigcup_{i=1}^\infty\bar{A}_i$.
Show, by an example, that this inclusion can be proper.
------

(a) By Theorem 2.24(d) and 2.27(a), $\bigcup_{i=1}^n\bar{A}_i$ is closed. Also, it is clear that $B_n\subset\bigcup_{i=1}^n\bar{A}_i$. So by Theorem 2.27(c), $\bar{B}_n\subset\bigcup_{i=1}^n\bar{A}_i$. Conversely, since $A_i\subset B_n$ for $i=1,2,\ldots,n$, we have $\bar{A}_i\subset\bar{B}_n$ for $i=1,2,\ldots,n$. Thus $\bigcup_{i=1}^n\bar{A}_i\subset\bar{B}_n$.

\smallskip
\noindent\textbf{Solution.}\quad
(b) To show this problem, use the same technique as in (a).

\smallskip
\noindent\textbf{Solution.}\quad
Finally, we give an example. Let $A_i=\left[\tfrac{1}{i+1},1\right]$ for $i=1,2,\ldots$. Then
\begin{align*}
\bigcup_{i=1}^\infty\bar{A}_i=\bigcup_{i=1}^\infty A_i=(0,1]
\end{align*}
On the other hand,
\begin{align*}
\bar{B}=\overline{\bigcup_{i=1}^\infty A_i}=\overline{(0,1]}=[0,1]
\end{align*}
Thus $\bigcup_{i=1}^\infty\bar{A}_i$ is a proper subset of $\bar{B}$.



\bigskip\noindent\hrule\bigskip

Is every point of every open set $E\subset \mathbb{R}^2$ a limit point of a set $E$? Answer the same question for closed sets in $\mathbb{R}^2$.

\smallskip
\noindent\textbf{Solution.}\quad

(i) Yes. To show this, let $E\subset \mathbb{R}^2$ be open, and $\mathbf{x}\in E$. Then by definition, there exists $r>0$ such that $\mathbf{y}\in E$ whenever $\left\vert\mathbf{x}-\mathbf{y}\right\vert<r$. Now, given $\varepsilon>0$, and write $\mathbf{x}=(x_1,x_2)$. Choose $\mathbf{z}=\left(x_1+\tfrac{1}{2}\min\{r,\varepsilon\},x_2\right)$, then clearly $\mathbf{z}\ne\mathbf{x}$ because $\frac{1}{2}\min\{r,\varepsilon\}>0$. Moreover, since
\begin{align*}
\left\vert\mathbf{z}-\mathbf{x}\right\vert=\frac{1}{2}\min\{r,\varepsilon\}\le\frac{r}{2}<r
\end{align*}
we see that $\mathbf{z}\in E$. Similarly, $\left\vert\mathbf{z}-\mathbf{x}\right\vert<\varepsilon$. Since $\varepsilon$ is arbitrarily chosen, we conclude that $\mathbf{x}$ is a limit point of $E$.
(ii) No. For example, consider any finite subset in $\mathbb{R}^2$.



\bigskip\noindent\hrule\bigskip

Let $E^\circ$ denote the set of all interior points of a set $E$. [See Definition 2.18(e); $E^\circ$ is called the interior of $E$.]
(a) Prove that $E^\circ$ is always open.
(b) Prove that $E$ is open if and only if $E^\circ=E$.
(c) If $G\subset E$ and $G$ is open, prove that $G\subset E^\circ$.
(d) Prove that the complement of $E^\circ$ is the closure of the complement of $E$.
(e) Do $E$ and $\bar{E}$ always have the same interiors?
(f) Do $E$ and $E^\circ$ always have the same closures?

\smallskip
\noindent\textbf{Solution.}\quad

(a) Let $x\in E^\circ$, i.e., $x$ is an interior point of $E$. Then there exists a neighborhood $N$ of $x$ such that $N\subset E$. It suffices to show that $N\subset E^\circ$. Since $N$ is open (by Theorem 2.19), every point of $N$ is an interior point of $N$, and so is an interior point of $E$, since $N\subset E$. This completes the proof.

\smallskip
\noindent\textbf{Solution.}\quad
(b) If $E$ is open, then by definition, every point of $E$ is an interior point of $E$. Thus $E\subset E^\circ$. Also, the inclusion $E^\circ\subset E$ is immediate [see Definition 2.18(e), the definition of interior points]. Hence $E^\circ=E$. Conversely, if $E^\circ=E$, then by part (a), $E$ is open.

\smallskip
\noindent\textbf{Solution.}\quad
(c) If $G\subset E$ and $G$ is open, then given $x\in G$, we see that $x$ is an interior point of $G$, and so is an interior point of $E$, since $G\subset E$. Hence $G\subset E^\circ$.

\smallskip
\noindent\textbf{Solution.}\quad
(d) Since $E^\circ\subset E$, we have $E^c\subset(E^\circ)^c$. Also, by part (a) and Theorem 2.23, $(E^\circ)^c$ is closed. Thus by Theorem 2.27(c), $\overline{E^c}\subset(E^\circ)^c$. Conversely, let $x\in(E^\circ)^c$, then we only have to consider the following two cases.
Case 1: $x\in E^c$. Then clearly $x\in\overline{E^c}$, and the claim follows.
Case 2. $x\in E$. Since $x$ is not an interior point of $E$, for every neighborhood $N$ of $x$, there always exists a point $y\in E^c$ such that $y\in N$. Since $x\in E$, we ensure that $x\ne y$. It follows that $x$ is a limit point of $E^c$, and then $x\in\overline{E^c}$.
Hence $\overline{E^c}=(E^\circ)^c$.

\smallskip
\noindent\textbf{Solution.}\quad
(e) No. For example, let $E=(0,1)\cup(1,2)$ in $\mathbb{R}$.

\smallskip
\noindent\textbf{Solution.}\quad
(f) No. For example, consider any finite set in $\mathbb{R}$.



\bigskip\noindent\hrule\bigskip

Let $X$ be an infinite set. For $p\in X$ and $q\in X$, define
\begin{align*}
d(p,q)=\left\{
\begin{array}{ll}
1&(\mbox{if $p\ne q$})\vspace{0.08in}\\
0&(\mbox{if $p=q$})
\end{array}\right.
\end{align*}
Prove that this is a metric. Which subsets of the resulting metric space are open? Which are closed? Which are compact?

\smallskip
\noindent\textbf{Solution.}\quad

For the first question, we check Definition 2.15 as below.
(a) Clearly $d(p,q)=1>0$ if $p\ne q$; and $d(p,p)=0$.
(b) If $p\ne q$, then $d(p,q)=1=d(q,p)$. If $p=q$, then $d(p,q)=0=d(q,p)$.
(c) Given $r\in X$. If $p\ne q$, then either $r\ne p$ or $r\ne q$, i.e., either $d(p,r)=1$ or $d(r,q)=1$. So
\begin{align*}
d(p,q)=1\le d(p,r)+d(r,q)
\end{align*}
If $p=q$, then
\begin{align*}
d(p,q)=0\le d(p,r)+d(r,q)
\end{align*}
is immediate.
Hence such $d$ is a metric.



\bigskip\noindent\hrule\bigskip

Exercise 2.11 is skipped.


\bigskip\noindent\hrule\bigskip

Let $K\subset \mathbb{R}$ consist of $0$ and the numbers $1/n$, for $n=1,2,3,\ldots$. Prove that $K$ is compact directly from the definition (without using the Heine-Borel theorem).

\smallskip
\noindent\textbf{Solution.}\quad

Let $\{G_\alpha\}$ be an open cover of $K$. For $n=1,2,\ldots$, choose an index $\alpha_n$ such that $\frac{1}{n}\in G_{\alpha_n}$, and an index $\alpha_0$ such that $0\in G_{\alpha_0}$. By the definition of open sets, there exists $r>0$ such that $y\in G_{\alpha_0}$ whenever $\left\vert y\right\vert=\left\vert y-0\right\vert<r$. Also, by archimedean, there exists a positive number $N$ such that $\frac{1}{N}<r$. So we have $\frac{1}{n}\in G_{\alpha_0}$ for all $n\ge N$, and hence there is a finitely many subcollection
\begin{align*}
\{G_{\alpha_0},G_{\alpha_1},G_{\alpha_2},\ldots,G_{\alpha_{N-1}}\}
\end{align*}
of $\{G_\alpha\}$ that covers $K$.



\bigskip\noindent\hrule\bigskip

Construct a compact set of real numbers whose limit points form a countable set.

\smallskip
\noindent\textbf{Solution.}\quad

Let
\begin{align*}
K=\{0\}\cup\bigcup_{n=1}^\infty\left(\left\{2^{-n}\right\}\cup\left\{2^{-n}+2^{-m}:m=n+1,n+2,\ldots\right\}\right)
\end{align*}
Then clearly $0$ and $2^{-n}$ ($n=1,2,\ldots$) are all the limit points of $K$. So it suffices to show that $K$ is compact. Let $\{G_\alpha\}$ be an open cover of $K$. For $n=1,2,\ldots,$ choose an index $\alpha_n$ such that $2^{-n}\in G_{\alpha_n}$, and an index $\alpha_0$ such that $0\in G_{\alpha_0}$. Then there exist $r_0>0$ and $r_n>0$ ($n=1,2,\ldots$), such that $y\in G_{\alpha_0}$ and $z\in G_{\alpha_n}$, whenever $\left\vert y\right\vert=\left\vert y-0\right\vert<r_0$ and $\left\vert z-2^{-n}\right\vert<r_n$. Also, there exists a positive integers $N_0$ and $N_n$ such that $2^{-N_0}<r_0$ and $2^{-n-N_n}<r_n$. It follows that for $j\ge N_0$ and $k\ge N_n$ ($n=1,2,\ldots$),
\begin{align*}
2^{-j}\in G_{\alpha_0},\quad 2^{-n}+2^{-n-k}\in G_{\alpha_n}
\end{align*}
Finally, for $n=1,2,\ldots,N_0$, denote the indices $\alpha_{n,1},\alpha_{n,2},\ldots,\alpha_{n,N_n-1}$ such that
\begin{align*}
2^{-n}+2^{-n-i}\in G_{\alpha_{n,i}}\quad\mbox{for }i=1,2,\ldots,N_n-1
\end{align*}
Then we conclude that
\begin{align*}
\{G_{\alpha_0}\}\cup\bigcup_{n=1}^{N_0}\left(\{G_{\alpha_n}\}\cup\{G_{\alpha_{n,i}}:i=1,2,\ldots,N_n-1\}\right)
\end{align*}
is a finite subcollection of $\{G_\alpha\}$ that covers $K$. Hence $K$ is compact.



\bigskip\noindent\hrule\bigskip

Give an example of an open cover of the segment $(0,1)$ which has no finite subcover.

\smallskip
\noindent\textbf{Solution.}\quad

Let $\left\{\left(\tfrac{1}{n+1},1\right)\right\}_{n=1}^\infty$.



\bigskip\noindent\hrule\bigskip

Show that Theorem 2.36 and its Corollary becomes false (in $\mathbb{R}$, for example) if the word "compact" is replaced by "closed" or by "bounded."

\smallskip
\noindent\textbf{Solution.}\quad

(i) If only the assumption "closed" holds, let $C_n=[n,\infty)$ for $n=1,2,\ldots$. Then each $C_n$ is closed, and the intersection of every finite subcollection of $\{C_n\}$ is nonempty, and also $C_1\supset C_2\supset C_3\supset\cdots$. But $\bigcap_{n=1}^\infty C_n=\varnothing$.
(ii) If only the assumption "bounded" holds, let $B_n=\left(0,n^{-1}\right)$ for $n=1,2,\ldots$. Then each $B_n$ is bounded, and the intersection of every finite subcollection of $\{B_n\}$ is nonempty, and also $B_1\supset B_2\supset B_3\supset \cdots$. But $\bigcap_{n=1}^\infty B_n=\varnothing$.



\bigskip\noindent\hrule\bigskip

Regard $Q$, the set of all rational numbers, as a metric space, with $d(p,q)=\left\vert p-q\right\vert$. Let $E$ be the set of all $p\in Q$ such that $2<p^2<3$. Show that $E$ is closed and bounded in $Q$, but that $E$ is not compact. Is $E$ open in $Q$?

\smallskip
\noindent\textbf{Solution.}\quad

(i) To show that $E$ is closed in $Q$, it suffices to show that $Q-E$ is open in $Q$. Let $p\in Q-E$, then $p\in Q$ and
\begin{align*}
p\in\left(-\infty,-\sqrt{3}\right]\cup\left[-\sqrt{2},\sqrt{2}\right]\cup\left[\sqrt{3},\infty\right)
\end{align*}
Since $\pm\sqrt{2}\ne Q$ and $\pm\sqrt{3}\ne Q$, we have
\begin{align*}
p\in\left(-\infty,-\sqrt{3}\right)\cup\left(-\sqrt{2},\sqrt{2}\right)\cup\left(\sqrt{3},\infty\right)
\end{align*}
Let $r=\min\left\{\left\vert p\pm\sqrt{2}\right\vert,\left\vert p\pm\sqrt{3}\right\vert\right\}$, then for all $q\in Q$ with $\left\vert p-q\right\vert<r$,
\begin{align*}
q\in\left(-\infty,-\sqrt{3}\right)\cup\left(-\sqrt{2},\sqrt{2}\right)\cup\left(\sqrt{3},\infty\right)
\end{align*}
i.e., $q\in Q-E$, and hence $Q-E$ is open.
(ii) Since for every $p\in E$, we have $p^2<3$ or $-\sqrt{3}<p<\sqrt{3}$. Hence $E\subset\left(-\sqrt{3},\sqrt{3}\right)$, or $E$ is bounded.
(iii) To show that $E$ is not compact, let $G_n=\left\{p\in Q:2<p^2<3-\tfrac{1}{n}\right\}$ for $n=2,3,\ldots$. Then the argument that each $G_n$ is open will be shown in the next part. It follows that $\{G_n\}$ is an open cover of $E$, but it has no finite subcover.
(iv) The set $E$ is open in $Q$. To show this, let $p\in E$, put $r=\min\left\{\left\vert p\pm\sqrt{2}\right\vert,\left\vert p\pm\sqrt{3}\right\vert\right\}$. Then for all $q\in Q$ with $\left\vert q-p\right\vert<r$,
\begin{align*}
\left\vert q\right\vert
=\left\vert (q-p)+p\right\vert
\le\left\vert q-p\right\vert+\left\vert p\right\vert
<r+\left\vert p\right\vert
\end{align*}
If $p\ge 0$, then since $p-\sqrt{3}<0$ always holds,
\begin{align*}
r+\left\vert p\right\vert
=r+p
\le\left\vert p-\sqrt{3}\right\vert+p
=\sqrt{3}-p+p
=\sqrt{3}
\end{align*}
i.e., $q^2<3$. If $p<0$, then since $p+\sqrt{3}>0$ always holds,
\begin{align*}
r+\left\vert p\right\vert
=r-p
\le\left\vert p+\sqrt{3}\right\vert-p
=p+\sqrt{3}-p
=\sqrt{3}
\end{align*}
i.e., $q^2<3$ again. On the other hand,
\begin{align*}
\left\vert q\right\vert
=\left\vert p+(q-p)\right\vert
\ge\left\vert p\right\vert-\left\vert q-p\right\vert
>\left\vert p\right\vert-r
\end{align*}
If $p\ge 0$, then since $p-\sqrt{2}>0$ always holds,
\begin{align*}
\left\vert p\right\vert-r
=p-r
\ge p-\left\vert p-\sqrt{2}\right\vert
=p-\left(p-\sqrt{2}\right)
=\sqrt{2}
\end{align*}
i.e., $q^2>2$. If $p<0$, then since $p+\sqrt{2}<0$ always holds,
\begin{align*}
\left\vert p\right\vert-r
=-p-r
\ge-p-\left\vert p+\sqrt{2}\right\vert
=-p+\left(p+\sqrt{2}\right)
=\sqrt{2}
\end{align*}
i.e., $q^2>2$ again. Hence we conclude that $q\in E$, i.e., $E$ is open in $Q$.



\bigskip\noindent\hrule\bigskip

Let $E$ be the set of all $x\in[0,1]$ whose decimal expansion contains only the digits $4$ and $7$. Is $E$ countable? Is $E$ dense in $[0,1]$? Is $E$ compact? Is $E$ perfect?

\smallskip
\noindent\textbf{Solution.}\quad

(i) The set $E$ is uncountable. To show this, suppose $A\subset E$ is a countable subset, let $A=\{a_1,a_2,\ldots\}$. We construct a number $a\in[0,1]$ whose $n$th digit is $4$ if the $n$th digit of $a_n$ is $7$, and is $7$ if the $n$th digit of $a_n$ is $4$. Then clearly $a\in E$ but $a\notin A$. It follows that every countable subset of $E$ is a proper subset, and hence $E$ is uncountable.
(ii) The set $E$ is not dense in $[0,1]$, since $E\subset\left[\tfrac{4}{9},\tfrac{7}{9}\right]$.
(iii) The set $E$ is compact. To show this, since $E$ is bounded, by the Heine-Borel theorem, it remains to show that $E$ is closed. Let $x\in E^c$ and $x\in[0,1]$, then the decimal expression of $x$ contains at least one digit different from $4$ and $7$. Write $x=0.x_1x_2x_3\ldots$, where $x_n\in \{0,1,2,\ldots,9\}$ for $n=1,2,\ldots$. Let $N$ be the first digit such that $x_N\notin \{4,7\}$. Given $a\in E$, write $a=0.a_1a_2a_3\ldots$, where $a_n\in\{4,7\}$ for $n=1,2,\ldots$. Suppose $M$ is the first digit such that $x$ and $a$ differ at the $M$th digit. Then $M\le N$ and
\begin{align*}
\left\vert x-a\right\vert
&=\left\vert \sum_{n=M}^\infty(x_n-a_n)\cdot 10^{-n}\right\vert \\
&\ge\left\vert x_M-a_M\right\vert10^{-M}-\sum_{n=M+1}^\infty\left\vert x_n-a_n\right\vert 10^{-n} \\
&\ge 10^{-M}-7\sum_{n=M+1}^\infty 10^{-n} \\
&=\frac{2\cdot 10^{-M}}{9} \\
&\ge\frac{2\cdot 10^{-N}}{9}
\end{align*}
where the third line of the above inequality follows by $\left\vert x_M-a_M\right\vert\ge 1$ and $\left\vert x_n-a_n\right\vert\le 7$ for all $n\ge M+1$. Since $N$ is fixed (depends only on $x$), we see that $x$ is an interior point of $E^c$. Hence $E^c$ is open, or $E$ is closed.
(iv) The set $E$ is perfect. To show this, let $x\in E$. Since for $r>0$, there exists a positive integer $n$ such that $10^{-n}<r$. So if we construct the number $y$ by changing the $(n+1)$th digit of $x$ from $4$ to $7$, or from $7$ to $4$, then $y\in E$, $y\ne x$, and
\begin{align*}
\left\vert x-y\right\vert=(7-4)\cdot 10^{-(n+1)}
<10^{-n}
<r
\end{align*}
Thus $x\in E'$. Finally, note that $E$ is closed [by part (iii)], so $E$ is perfect.



\bigskip\noindent\hrule\bigskip

Is there a nonempty perfect set in $\mathbb{R}$ which contains no rational number?

\smallskip
\noindent\textbf{Solution.}\quad

Yes. Let $\{r_1,r_2,\ldots\}$ be the set of all rational numbers in $E_0=[\pi,\pi+1]$, construct a subset of $E_0$ as below. Having chosen $E_n$ ($n\ge 0$) the subsets of $E_0$ such that $E_n$ is a disjoint union of at most $2^{n+1}-1$ intervals with irrational endpoints, each of positive length at most $3^{-n}$ and $E_n$ does not contain $r_j$ for $1\le j\le n$. Define the set $F_{n+1}$ obtained from $E_n$ by removing the middle thirds of all intervals of $E_n$. Thus $F_{n+1}$ is a disjoint union of at most $2^{n+2}-2$ intervals with irrational endpoints, each of positive length at most $3^{-(n+1)}$. If $r_{n+1}\notin F_{n+1}$, take $E_{n+1}=F_{n+1}$. If $r_{n+1}\in F_{n+1}$, then $r_{n+1}\in [a,b]$ for some interval $[a,b]\subset F_{n+1}$. Since $a$ and $b$ are irrational, we have $r_{n+1}\in(a,b)$. Let $\delta>0$ be an irrational number such that $\delta<\min\{r_{n+1}-a,b-r_{n+1}\}$, and put
\begin{align*}
E_{n+1}=F_{n+1}-(r_{n+1}-\delta,r_{n+1}+\delta)
\end{align*}
By doing this, we see that $E_{n+1}$ is a disjoint union of at most $2^{n+2}-1$ intervals with irrational endpoints, each of positive length at most $3^{-(n+1)}$ and $E_{n+1}$ does not contain $r_j$ for $1\le j\le n+1$. Hence the resulting set $P=\bigcap_{n=0}^\infty E_n$ contains no rational number.
Since $E_n$ is closed for each $n\ge 0$, by Theorem 2.24(b), $P$ is closed. Moreover, since $E_0\supset E_1\supset E_2\supset \cdots$, each $E_n$ is bounded (by $E_0$), i.e., each $E_n$ is actually a (nonempty) compact set. By Theorem 2.36, $P\ne\varnothing$. Finally, to show that $P$ is perfect, let $x\in P$, then $x\in E_n$ for each positive integer $n$. Let $x\in[a,b]$ for some interval $[a,b]\subset E_n$, note that $b-a\le 3^{-n}$. If $x\ne a$, then put $y=a$; If $x=a$, then put $y=b$. In both cases, we see that $y\in P$, $y\ne x$, and
\begin{align*}
\left\vert y-x\right\vert\le b-a\le 3^{-n}<2\cdot 3^{-n}
\end{align*}
It follows that $x$ is a limit point of $P$, and hence $P$ is perfect.



\bigskip\noindent\hrule\bigskip

(a) If $A$ and $B$ are disjoint closed sets in some metric space, prove that they are separated.
(b) Prove the same for disjoint open sets.
(c) Fix $p\in X$, $\delta>0$, define $A$ to be the set of all $q\in X$ for which $d(p,q)<\delta$, define $B$ similarly, with $>$ in place of $<$. Prove that $A$ and $B$ are separated.
(d) Prove that every connected metric space with at least two points is uncountable. \textit{Hint:} Use (c).

\smallskip
\noindent\textbf{Solution.}\quad

(a) Since $A$ and $B$ are closed sets, we have $\bar{A}=A$ and $\bar{B}=B$. Thus $A\cap B=\varnothing$ implies
\begin{align*}
\bar{A}\cap B=A\cap B=\varnothing,\quad
A\cap \bar{B}=A\cap B=\varnothing
\end{align*}
i.e., $A$ and $B$ are separated.

\smallskip
\noindent\textbf{Solution.}\quad
(b) Let $A$ and $B$ be two disjoint open sets, and let $a\in A$. By definition, there exists a neighborhood $N$ of $x$ such that $N\subset A$. It follows that $N\cap B=\varnothing$, and thus $x\notin B'$. So $A\cap\bar{B}=\varnothing$. Similarly, we also have $\bar{A}\cap B=\varnothing$. Hence $A$ and $B$ are separated.

\smallskip
\noindent\textbf{Solution.}\quad
(c) Note that $A$ is open (by Theorem 2.19), thus by part (b), it suffices to show that $B=\{q\in X:d(p,q)>\delta\}$ is open. Let $q\in B$, and denote $h=d(p,q)-\delta>0$. Then for every $r\in X$ with $d(q,r)<h$,
\begin{align*}
d(p,r)
\ge d(p,q)-d(q,r)
>d(p,q)-h
=\delta
\end{align*}
It follows that $r\in B$, and hence $B$ is open.

\smallskip
\noindent\textbf{Solution.}\quad
(d) Let $E\subset X$ be a connected set, and $p,q\in E$ be two distinct points. Denote $d=d(p,q)>0$, we claim that, for each $\delta\in(0,d)$, there exists $r\in E$ such that $d(p,r)=\delta$.
\textit{Proof of the claim.} If otherwise, there exists $\delta_0\in(0,d)$ such that, for all $r\in E$, either $d(p,r)<\delta_0$ or $d(p,r)>\delta_0$. If we let $A=E\cap\{r\in X:d(p,r)<\delta_0\}$ and $B=E\cap\{r\in X:d(p,r)>\delta_0\}$, then $E=A\cup B$ whereas $A$ and $B$ are separated [by part (c)]. This contradicts the assumption that $E$ is connected. $\Box$
Hence $E$ can be mapped onto the segment $(0,d)$, and so $E$ is uncountable.



\bigskip\noindent\hrule\bigskip

Are closures and interiors of connected sets always connected? (Look at subsets of $\mathbb{R}^2$.)

\smallskip
\noindent\textbf{Solution.}\quad

(i) The closure of a connected set is connected. To show this, let $E\subset X$ be a set. Suppose $\bar{E}$ is not connected, then we can write $\bar{E}=A\cup B$ for some separated sets $A$ and $B$. Since $E\subset\bar{E}$, we have $E\subset A\cup B$, and then
\begin{align*}
E=E\cap(A\cup B)=(E\cap A)\cup(E\cap B)
\end{align*}
Since $A$ and $B$ are separated, we have
\begin{align*}
\overline{(E\cap A)}\cap(E\cap B)&\subset\bar{A}\cap B=\varnothing \\
(E\cap A)\cap\overline{(E\cap B)}&\subset A\cap\bar{B}=\varnothing
\end{align*}
i.e., $\overline{(E\cap A)}\cap(E\cap B)=(E\cap A)\cap\overline{(E\cap B)}=\varnothing$. So we conclude that $E\cap A$ and $E\cap B$ are separated, and hence $E$ is not connected.
(ii) The interior of a connected set is not always connected. For example, in $\mathbb{R}^2$, consider $D_1\cup D_2$, where $D_1=\{(x,y):x^2+y^2\le 1\}$ and $D_2=\{(x,y):(x-2)^2+y^2\le 1\}$.



\bigskip\noindent\hrule\bigskip

Let $A$ and $B$ be separated subsets of some $\mathbb{R}^k$, suppose $\mathbf{a}\in A$, $\mathbf{b}\in B$, and define
\begin{align*}
\mathbf{p}(t)
=(1-t)\mathbf{a}
+t\mathbf{b}
\end{align*}
for $t\in \mathbb{R}$. Put $A_0=\mathbf{p}^{-1}(A)$, $B_0=\mathbf{p}^{-1}(B)$. [Thus $t\in A_0$ if and only if $\mathbf{p}(t)\in A$.]
(a) Prove that $A_0$ and $B_0$ are separated subsets of $\mathbb{R}$.
(b) Prove that there exists $t_0\in(0,1)$ such that $\mathbf{p}(t_0)\notin A\cup B$.
(c) Prove that  every convex subset of $\mathbb{R}^k$ is connected.

\smallskip
\noindent\textbf{Solution.}\quad

(a) It suffices to show that $\bar{A}_0\cap B_0=\varnothing$ (the same argument as $A_0\cap\bar{B}_0=\varnothing$, by symmetry). Suppose $t\in\bar{A}_0\cap B_0$, then either $t\in A_0\cap B_0$ or $t\in A_0'\cap B_0$, since 
\begin{align*}
\bar{A}_0\cap B_0
=(A_0\cup A_0')\cap B_0
=(A_0\cap B_0)\cup(A_0'\cap B_0)
\end{align*}
So we consider the following two cases.
Case 1: If $t\in A_0\cap B_0$, then $\mathbf{p}(t)\in A\cap B$. But since $A$ and $B$ are separated, we must have $A\cap B=\varnothing$. Then such $\mathbf{p}(t)$ does not exist, and so does $t$.
Case 2: If $t\in A_0'\cap B_0$, then $t\in B_0$ and given $\varepsilon>0$, there exists $s\in A_0$ such that $\left\vert t-s\right\vert<\varepsilon$. So we have
\begin{align*}
\left\vert\mathbf{p}(t)-\mathbf{p}(s)\right\vert
\le\left\vert t-s\right\vert\left(\left\vert\mathbf{a}\right\vert+\left\vert\mathbf{b}\right\vert\right)
<\varepsilon\left(\left\vert\mathbf{a}\right\vert+\left\vert\mathbf{b}\right\vert\right)
\end{align*}
Since $\mathbf{p}(s)\in A$ and $\varepsilon$ can be arbitrarily small, we see that $\mathbf{p}(t)\in A'$. Thus
\begin{align*}
\mathbf{p}(t)
\in A'\cap B
\subset\bar{A}\cap B
\end{align*}
Since $A$ and $B$ are separated, we have $\bar{A}\cap B=\varnothing$. Then such $\mathbf{p}(t)$ does not exist, and so does $t$.
In both cases, we conclude that $\bar{A}_0\cap B_0=\varnothing$.

\smallskip
\noindent\textbf{Solution.}\quad
(b) By part (a), the sets $A_0$ and $B_0$ are separated. Also, it is clear that $0\in A_0$ and $1\in B_0$. So by the proof of Theorem 2.47, there exists $t_0$ with $0<t_0<1$, such that $t_0\notin A_0\cup B_0$, or $\mathbf{p}(t_0)\notin A\cup B$.

\smallskip
\noindent\textbf{Solution.}\quad
(c) Let $C\subset \mathbb{R}^k$ be a convex subset, then for all $\mathbf{a},\mathbf{b}\in C$ and for all $t\in[0,1]$, we have $\mathbf{p}(t)\in C$. So, even if we write $C=A\cup B$, and put $\mathbf{a}\in A$ and $\mathbf{b}\in B$, the statement of part (b) does not hold. It follows that such $A$ and $B$ are not separated, and hence $C$ is connected.



\bigskip\noindent\hrule\bigskip

A metric space is called \textit{separated} if it contains a countable dense subset. Show that $\mathbb{R}^k$ is separable. \textit{Hint:} Consider the set of points which have only rational coordinates.

\smallskip
\noindent\textbf{Solution.}\quad

Let
\begin{align*}
E=\{(r_1,r_2,\ldots,r_k): r_1,r_2,\ldots,r_k\mbox{ are rationals}\}
\end{align*}
then clearly $E\subset \mathbb{R}^k$ is countable. Now, let $(x_1,x_2,\ldots,x_k)\in \mathbb{R}^k$ and given $\varepsilon>0$, then for $j=1,2\ldots,k$, there exists $r_j^\ast\in Q$ such that $r_j^\ast\in\left(x_j-\tfrac{\varepsilon}{\sqrt{k}},x_j+\tfrac{\varepsilon}{\sqrt{k}}\right)$. It follows that
\begin{align*}
\sum_{j=1}^k\left(x_j-r_j^\ast\right)^2
<k\cdot\left(\frac{\varepsilon}{\sqrt{k}}\right)^2
=\varepsilon^2
\end{align*}
and so we obtain a point $(r_1^\ast,r_2^\ast,\ldots,r_k^\ast)\in E$ satisfying $(r_1^\ast,r_2^\ast,\ldots,r_k^\ast)\in N_\varepsilon(x_1,x_2,\ldots,x_k)$. Hence $E$ is dense in $\mathbb{R}^k$, or $\mathbb{R}^k$ is separable.



\bigskip\noindent\hrule\bigskip

A collection $\{V_\alpha\}$ of open subsets of $X$ is said to be a \textit{base} for $X$ if the following is true: For every $x\in X$ and every open set $G\subset X$ such that $x\in G$, we have $x\in V_\alpha\subset G$ for some $\alpha$. In other words, every open set in $X$ is the union of a subcollection of $\{V_\alpha\}$.
Prove that every separable metric space has a \textit{countable} base. \textit{Hint:} Take all neighborhoods with rational radius and center in some countable dense subset of $X$.

\smallskip
\noindent\textbf{Solution.}\quad

Let $X$ be a separable metric space, then by definition, we may let $E$ be a countable dense subset of $X$, and write $E=\{e_1,e_2,\ldots\}$. For each positive integer $m$, define $V_n^m=N_{\tfrac{1}{n}}(e_m)$ for $n=1,2,\ldots$. Then it is clear that the collection $\{V_n^m\}_{m,n}$ is countable. Now, we show that $\{V_n^m\}_{m,n}$ is a base for $X$. Let $x\in X$, and $G\subset X$ be an open set such that $x\in G$, then by the definition of open sets, there exists $r>0$ such that $N_r(x)\subset G$. Also, by archimedean, there exists a positive integer $N$ such that $\frac{1}{N}<r$. So $N_{\tfrac{1}{N}}(x)\subset G$ holds. By the definition of dense sets, there exists $e_M\in E$, for some integer $M$, such that $e_M\in N_{\tfrac{1}{2N}}(x)$. So does $x\in N_{\tfrac{1}{2N}}(e_M)=V_{2N}^M$. Finally, given $y\in V_{2N}^M$, we have
\begin{align*}
d(x,y)
\le d(x,e_M)+d(e_M,y)
<\frac{1}{2N}+\frac{1}{2N}
=\frac{1}{N}
\end{align*}
or $y\in N_{\tfrac{1}{N}}(x)$. It follows that $V_{2N}^M\subset N_{\tfrac{1}{N}}(x)$, and thus $x\in V_{2N}^M\subset G$. Hence $\{V_n^m\}_{m,n}$ is a base for $X$.



\bigskip\noindent\hrule\bigskip

Let $X$ be a metric space in which every infinite subset has a limit point. Prove that $X$ is separable. \textit{Hint:} Fix $\delta>0$, and pick $x_1\in X$. Having chosen $x_1,\ldots,x_j\in X$, choose $x_{j+1}\in X$, if possible, so that $d(x_i,x_{j+1})\ge\delta$ for $i=1,\ldots,j$. Show that this process must stop after a finite number of steps, and that $X$ can therefore be covered by finitely many neighborhood of radius $\delta$. Take $\delta=1/n$ ($n=1,2,3,\ldots$), and consider the centers of the corresponding neighborhoods.

\smallskip
\noindent\textbf{Solution.}\quad

Let $(X,d)$ be an infinite metric space satisfying for each infinite subset $E\subset X$, $E$ has a limit point. Following the hint, let $x_1\in X$ and given $\delta>0$. Having chosen $x_1,x_2,\ldots,x_j\in X$, choose $x_{j+1}\in X$ if possible, so that $d(x_i,x_{j+1})\ge\delta$ for $i=1,2,\ldots,j$. We claim that this process will stop at a finite step.
\textit{Proof of the claim.} Since if not, we may create an infinite subset $E=\{x_1,x_2,\ldots\}$ of $X$, so that $d(x_i,x_j)\ge\delta$ for all integers $i,j$ with $i\ne j$. Then for every $x\in X$ with $d(x_i,x)<\frac{\delta}{2}$ for some $i$, we have
\begin{align*}
d(x_j,x)\ge d(x_j,x_i)-d(x_i,x)\ge\delta-\frac{\delta}{2}=\frac{\delta}{2}
\end{align*}
for all $j$ with $j\ne i$. It follows that the neighborhood $N_{\tfrac{\delta}{2}}(x)$ contains a finite number of points of $E$, and thus $E$ has no limit point, a contradiction. $\Box$
Now, take $\delta=\frac{1}{n}$ for $n=1,2,\ldots$. There is a corresponding finite subset $E_n$ satisfying the above condition. We show that $D=\bigcup_{n=1}^\infty E_n$ is a countable dense subset of $X$. First, $D$ is clearly countable. Next, to show that $D$ is dense in $X$, let $x\in X$ and given $\varepsilon>0$. If $x\in D$, then we are done. If $x\notin D$, by archimedean, there exists a positive integer $N$ so that $\frac{1}{N}<\varepsilon$. Also, there exists $y\in E_N$ with $y\ne x$, such that $d(x,y)<\frac{1}{N}<\varepsilon$. For otherwise, by the construction of $E_N$, we may regard $x$ as a member of $E_N$ (continue the process of the construction of $E_N$), and then $x\in D$, contradicting the assumption $x\notin D$. So in this case, $x$ is a limit point of $D$. Hence $D$ is dense in $X$.



\bigskip\noindent\hrule\bigskip

Prove that every compact metric space $K$ has a countable base, and that $K$ is therefore separable. \textit{Hint:} For every positive integer $n$, there are finitely many neighborhoods of radius $1/n$ whose union covers $K$.

\smallskip
\noindent\textbf{Solution.}\quad

(i) For $n=1,2,\ldots$, the collection $\left\{N_{\tfrac{1}{n}}(x)\right\}_{x\in K}$ is clearly an open cover of $K$. Since $K$ is compact, there are finitely many points $x_{n,1},x_{n,2},\ldots,x_{n,k_n}\in K$ such that the finite subcollection $\mathcal{B}_n=\left\{N_{\tfrac{1}{n}}(x_{n,j})\right\}_{j=1}^{k_n}$ also covers $K$. Let 
\begin{align*}
\mathcal{B}
=\bigcup_{n=1}^\infty\mathcal{B}_n
\end{align*}
We show that $\mathcal{B}$ is a countable base for $K$. First, $\mathcal{B}$ is clearly countable. Next, given $x\in K$ and an open subset $G\subset K$ with $x\in G$. By the definition of open sets, there exists $r>0$ such that $N_r(x)\subset G$. Also, by archimedean, there exists a positive integer $N$ such that $\frac{1}{N}<r$. Since $\mathcal{B}_{2N}$ covers $K$, we have $x\in N_{\tfrac{1}{2N}}(x_{2N,j})$ for some $1\le j\le k_{2N}$. Note that $N_{\tfrac{1}{2N}}(x_{2N,j})\subset N_{\tfrac{1}{N}}(x)$, since for $y\in N_{\tfrac{1}{2N}}(x_{2N,j})$, we have
\begin{align*}
d(x,y)\le d(x,x_{2N,j})+d(x_{2N,j},y)<\frac{1}{2N}+\frac{1}{2N}=\frac{1}{N}
\end{align*}
or $y\in N_{\tfrac{1}{N}}(x)$. Thus we conclude that $x\in N_{\tfrac{1}{2N}}(x_{2N,j})\subset N_r(x)\subset G$, and hence $\mathcal{B}$ is a base for $K$.
(ii) Let $x\in K$ and given $\varepsilon>0$, choose a positive $N$ such that $\frac{1}{N}<\varepsilon$. Since $\mathcal{B}_{N}$ covers $K$, we have $x\in N_{\tfrac{1}{N}}(x_{N,j})$ for some $1\le j\le k_{N}$, and thus 
\begin{align*}
x_{N,j}\in  N_{\tfrac{1}{N}}(x)\subset N_\varepsilon(x)
\end{align*}
It follows that $\bigcup_{n=1}^\infty\{x_{n,j}\}_{j=1}^{k_n}$ forms a countable dense subset of $K$, and hence $K$ is separable.
\textit{Alternative Solution of part (ii).} Since $K$ is compact, by Theorem 2.37, every infinite subset of $K$ has a limit point. Thus by Exercise 24, $K$ is separable.



\bigskip\noindent\hrule\bigskip

Let $X$ be a metric space in which every infinite subset has a limit point. Prove that $X$ is compact. \textit{Hint:} By Exercise 23 and 24, $X$ has a countable base. It follows that every open cover of $X$ has a \textit{countable} subcover $\{G_n\}$, $n=1,2,3,\ldots$. If no finite subcollection of $G_n$ covers $X$, then the complement $F_n$ of $G_1\cup\cdots\cup G_n$ is nonempty for each $n$, but $\bigcap F_n$ is empty. If $E$ is a set which contains a point from each $F_n$, consider a limit point of $E$, and obtain a contradiction.

\smallskip
\noindent\textbf{Solution.}\quad

Following the hint, let $E=\{x_n\}\subset X$, where $x_n\in F_n$ for each $n$. Then $E$ must be an infinite set. To show this, let $\{x_{i_1},x_{i_2},\ldots,x_{i_m}\}$ be a finite subset of $E$. Since $\{G_n\}$ covers $X$, we can choose $G_{j_1},G_{j_2},\ldots,G_{j_m}$ such that $x_{i_k}\in G_{j_k}$ for $1\le k\le m$. Put $M=\max\{j_1,j_2,\ldots,j_m\}$, then the point $x_M$ satisfies
\begin{align*}
x_M\in F_M=(G_1\cup G_2\cup\cdots\cup G_M)^c\subset(G_{j_1}\cup G_{j_2}\cup\cdots\cup G_{j_m})^c
\end{align*}
Thus $x_M\neq x_{i_k}$ for $1\le k\le m$. So every finite subset of $E$ is a proper subset, and hence $E$ is infinite. Now, by hypothesis, let $x$ be a limit point of $E$. Since $F_1\supset F_2\supset F_3\supset\cdots$, the set $\{x_{n+1},x_{n+2},\ldots\}\subset F_n$ for each $n$. Also, since $x$ is a limit point of $E$, $x$ is also a limit of $\{x_{n+1},x_{n+2},\ldots\}$, and so $x$ is a limit point of $F_n$, for each $n$. Since each $F_n$ is closed, we see that $x\in F_n$ for each $n$, i.e., $x\in\bigcap_{n=1}^\infty F_n$. This contradicts $\bigcap_{n=1}^\infty F_n=\varnothing$.



\bigskip\noindent\hrule\bigskip

Define a point $p$ in a metric space $X$ to be a \textit{condensation point} of a set $E\subset X$ if every neighborhood of $p$ contains uncountably many points of $E$.
Suppose $E\subset \mathbb{R}^k$, $E$ is uncountable, and let $P$ be the set of all condensation points of $E$. Prove that $P$ is perfect and that at most countably many points of $E$ are not in $P$. In other words, show that $P^c\cap E$ is at most countable. \textit{Hint:} Let $\{V_n\}$ be a countable base of $\mathbb{R}^k$, let $W$ be the union of those $V_n$ for which $E\cap V_n$ is at most countable, and show that $P=W^c$.

\smallskip
\noindent\textbf{Solution.}\quad

(i) By Exercise 23 and 24, we may let $\{V_n\}$ be a countable base for $\mathbb{R}^k$. Let $W$ be the union of those $V_n$ for which $E\cap V_n$ is at most countable, then we show that $P=W^c$. If $\mathbf{x}\in P$, and given a member of the base $V_n$, such that $\mathbf{x}\in V_n$. Since $V_n$ is open, there exists a neighborhood $N$ of $\mathbf{x}$ such that $\mathbf{x}\in N\subset V_n$. Since $\mathbf{x}\in P$, the set $E\cap N$ is uncountable, and so is $E\cap V_n$. It follows that $\mathbf{x}\notin V_n$ for all $n$ that satisfies $E\cap V_n$ is at most countable, and hence $\mathbf{x}\in W^c$. Conversely, if $\mathbf{x}\in W^c$, and given a neighborhood $N$ of $\mathbf{x}$. Then by the definition of bases, there exists some positive integer $N$ such that $\mathbf{x}\in V_N\subset N$. Since $\mathbf{x}\in W^c$, the set $E\cap V_N$ is uncountable, and so is $E\cap N$. Hence $\mathbf{x}\in P$.
Now, we show that $P$ is perfect. First, $P$ is closed since $P=W^c$, where $W$ is open. Next, to show that $P\subset P'$, suppose $\mathbf{p}\notin P'$. Then there exists a neighborhood $N$ of $\mathbf{p}$ such that for every $\mathbf{q}\in N$ with $\mathbf{q}\ne \mathbf{p}$, we have $\mathbf{q}\notin P$ or $\mathbf{q}\in W$. It follows that
\begin{align*}
N-\{\mathbf{p}\}\subset W\quad\mbox{or}\quad N\subset W\cup\{\mathbf{p}\}
\end{align*}
Since $E\cap W$ is at most countable, so is $E\cap\left(W\cup\{\mathbf{p}\}\right)$, and $E\cap N$ is also at most countable. Thus $\mathbf{p}\notin P$, and we conclude that $P$ is perfect.
(ii) Since $P=W^c$ and $E\cap W$ is at most countable, it is clear that $P^c\cap E$ is at most countable.
\textbf{Remark.} The statement is true for any separable metric space, see Exercise 28.



\bigskip\noindent\hrule\bigskip

Prove that every closed set in a separable metric space is the union of a (possibly empty) perfect set and a set which is at most countable. (\textit{Corollary:} Every countable closed set in $\mathbb{R}^k$ has isolated points.) \textit{Hint:} Use Exercise 27.

\smallskip
\noindent\textbf{Solution.}\quad

(i) Let $X$ be a separable metric space with countable base $\{V_n\}$, and let $F\subset X$ be a closed subset. If $F$ is at most countable, then we can write $F=\varnothing\cup F$ and this case is finished. If $F$ is uncountable, let $P$ be the set of all condensation points of $F$. Then clearly $P\subset F'$, and then $P\subset F$ because $F$ is closed. So we can write $F=P\cup(F-P)$. Also, by Exercise 27, $P$ is perfect. It therefore suffices to show that $F-P$ is at most countable. But this result is immediate because
\begin{align*}
F-P
=F\cap P^c
=F\cap W
\end{align*}
where $W$ is the union of those $V_n$ for which $F\cap V_n$ is at most countable (see Exercise 27).
(ii) To show the corollary, since every nonempty perfect set in $\mathbb{R}^k$ is uncountable (by Theorem 2.43), so every countable closed set $F$ in $\mathbb{R}^k$ must not be perfect. i.e., $F$ has isolated points.



\bigskip\noindent\hrule\bigskip

Prove that every open set in $\mathbb{R}$ is the union of an at most countable collection of disjoint segments. \textit{Hint:} Use Exercise 22.

\smallskip
\noindent\textbf{Solution.}\quad

Let $G\subset \mathbb{R}$ be an open set. For $x\in G$, define $A_x=\{z:[z,x]\subset G\}$, $B_x=\{z:[x,z]\subset G\}$, $a_x=\inf A_x$, $b_x=\sup B_x$, and $S_x=(a_x,b_x)$. Note that if $B_x$ is bounded above, then $b_x\in \mathbb{R}$ exists (by the least-upper-bound property). If otherwise, then $b_x=\infty$. A similar argument for $a_x$. Now, we claim that if $(a,b)\in \mathbb{R}$ is a segment so that $x\in(a,b)\subset G$, then $(a,b)\subset S_x$. i.e., $a_x\le a$ and $b\le b_x$.
\textit{Proof of the claim.} If not, consider the case $a_x>a$ (the other case $b>b_x$ is similar). Then there exists $a'$ so that $a_x>a'>a$. Moreover, since $(a,x]\subset G$, we have $[a',x]\subset G$, which implies $a'\in A_x$. This contradicts the fact that $a_x=\inf A_x$. $\Box$
By the definition of open sets, we can choose a $\delta>0$ such that $(x-\delta,x+\delta)\subset G$. Since $a_x\le x-\delta$ and $x+\delta\le b_x$ (by the above claim), we have $a_x<b_x$ and $x\in S_x$. To show that $S_x\subset G$, given $z\in S_x$, there exists $z_1\in A_x$ and $z_2\in B_x$ so that $z_1<z<z_2$. Thus
\begin{align*}
z\in[z_1,z_2]=[z_1,x]\cup[x,z_2]\subset G
\end{align*} 
Next, we claim that for $x,y\in G$ and $x\ne y$, either $S_x=S_y$ or $S_x\cap S_y=\varnothing$.
\textit{Proof of the claim.} If $S_x\cap S_y\ne\varnothing$, then $S_x\cup S_y$ is an open interval containing $x$ and $y$, such that $S_x\cup S_y\subset G$. By the above claim, $S_x\cup S_y\subset S_x$ and $S_x\cup S_y\subset S_y$, and hence $S_x=S_x\cup S_y=S_y$. $\Box$
Now, since $G=\bigcup_{x\in G}S_x$ and each $S_x$ contains a rational number, the number of distinct $S_x$ must be at most countable. This completes the proof.



\bigskip\noindent\hrule\bigskip

Imitate the proof of Theorem 2.43 to obtain the following result:
If $\mathbb{R}^k=\bigcup_{1}^\infty F_n$, where each $F_n$ is a closed subset of $\mathbb{R}^k$, then at least one $F_n$ has a nonempty interior.
\textit{Equivalent statement:} If $G_n$ is a dense open subset of $\mathbb{R}^k$, for $n=1,2,3,\ldots$, then $\bigcap_{1}^\infty G_n$ is not empty (in fact, it is dense in $\mathbb{R}^k$).
(This is a special case of Baire's theorem; see Exercise 22, Chap. 3, for the general case.)

\smallskip
\noindent\textbf{Solution.}\quad

First, we prove the second statement. Let $x\in \mathbb{R}^k$ and given $\varepsilon>0$, then since $G_n$ is a dense open subsets of $\mathbb{R}^k$ for $n=1,2,\ldots$, we can choose a $g_1\in G_1$ such that $g_1\in N_{\varepsilon}(x)$. Moreover, there exists $\delta_1>0$ such that $\overline{N_{\delta_1}(g_1)}\subset N_{\varepsilon}(x)\cap G_1$. Next, choose $g_2\in G_2$ such that $g_2\in N_{\delta_1}(g_1)$, and choose $\delta_2>0$ such that $\overline{N_{\delta_2}(g_2)}\subset N_{\delta_1}(g_1)\cap G_2$. Continue this process, the resulting set $\bigcap_{n=1}^\infty\overline{N_{\delta_n}(g_n)}$ contains a point $g\in \mathbb{R}^k$ (by the corollary to Theorem 2.36). Obviously $g\in G_n$ for $n=1,2,\ldots$, and then $g\in\bigcap_{n=1}^\infty G_n$. Also, $g\in\overline{N_{\delta_1}(g_1)}\subset N_\varepsilon(x)$. Thus $\bigcap_{n=1}^\infty G_n$ is dense in $\mathbb{R}^k$.
Next, we prove the equivalence of the two statements. Suppose the first statement holds, let $G_n$ be a dense open subset of $\mathbb{R}^k$, for $n=1,2,\ldots$. Denote the closed set $F_n=(G_n)^c$ for each $n$, then since $G_n$ is dense in $\mathbb{R}^k$, for $x\in F_n$ and for $\varepsilon>0$, we have $N_r(x)\cap G_n\ne\varnothing$. This means $F_n$ has an empty interior for each $n$. By the first statement and the De Morgan law,
\begin{align*}
\left(\bigcap_{n=1}^\infty G_n\right)^c=\bigcup_{n=1}^\infty F_n\ne \mathbb{R}^k
\end{align*}
i.e., $\bigcap_{n=1}^\infty G_n\ne\varnothing$. Conversely, suppose the second statement holds, let $F_n$ be a closed set for $n=1,2,\ldots$, such that $\mathbb{R}^k=\bigcup_{n=1}^\infty F_n$. Denote the open set $G_n=(F_n)^c$ for each $n$, then by the De Morgan law again,
\begin{align*}
\varnothing
=\left(\bigcup_{n=1}^\infty F_n\right)^c
=\bigcap_{n=1}^\infty G_n
\end{align*}
It follows by the second statement that there is an integer $N$ such that $G_N$ is not dense in $\mathbb{R}^k$, i.e., there exists $x\in \mathbb{R}^k$ and $r>0$, such that $N_r(x)\cap G_N=\varnothing$, or $N_r(x)\subset (G_N)^c=F_N$ for some $N$.


\end{document}
